\documentclass[12pt]{article}
\usepackage[T1]{fontenc}
\usepackage[utf8]{inputenc}
\usepackage{lmodern}
\usepackage{textcomp}
\usepackage{enumitem}
\usepackage{geometry}
\geometry{margin=1in}
\begin{document}
\begin{center}
Answer Key - Chemistry - Graduate Level Assessment 23 \
\small (Answers are ordered exactly as in the question paper.)
\end{center}

\section*{Multiple Choice}
\begin{enumerate}[label=\arabic*.]
\item \textbf{Answer:} C
\\ \textit{Options:} A) 200.618 g of CaBr\_2; B) 175.54 g of CaBr\_2; C) 105.32 g of CaBr\_2; D) 50.154 g of CaBr\_2; E) 150 g of CaBr\_2
\\ \textit{Note:} To prepare a 3.5 M solution of calcium bromide (CaBr₂), you need to use the formula mass = molarity × volume (in liters), which calculates to 3.5 moles/L × 0.150 L × 200.618 g/mole = 105.32 g of CaBr₂.
\item \textbf{Answer:} A
\\ \textit{Options:} A) 15; B) 40; C) 38; D) 13; E) 21
\\ \textit{Note:} The correct answer is 15 half-lives because each half-life reduces the amount of substance by half, so starting from 10.0 kg (10,000,000 mg) to reach 1.0 microgram (0.001 mg) requires 15 successive halvings, calculated as 10,000,000 mg divided by 2 raised to the power of 15.
\item \textbf{Answer:} A
\\ \textit{Options:} A) It causes the pH to fluctuate unpredictably.; B) It causes the pH to equalize with the water's pH.; C) It causes the pH to become extremely basic.; D) It neutralizes the pH completely.; E) It significantly decreases the pH.
\\ \textit{Note:} Doubling the volume of the buffer solution by adding water dilutes both the weak acid and its conjugate base equally, which can disrupt the buffer's capacity to maintain a stable pH, leading to unpredictable fluctuations.
\item \textbf{Answer:} E
\\ \textit{Options:} A) IR radiation can ionize organic molecules, providing structural information; B) absorption of IR radiation increases with molecular complexity; C) all molecular bonds absorb IR radiation; D) IR spectroscopy can identify isotopes within organic molecules; E) most organic functional groups absorb in a characteristic region of the IR spectrum
\\ \textit{Note:} The correct answer is accurate because each organic functional group has specific vibrational modes that correspond to distinct absorption frequencies in the infrared spectrum, enabling their identification and structural analysis.
\item \textbf{Answer:} A
\\ \textit{Options:} A) 5.9 $\mathrm{~cm}^3$; B) 12.1 $\mathrm{~cm}^3$; C) 7.3 $\mathrm{~cm}^3$; D) 7.9 $\mathrm{~cm}^3$; E) 9.5 $\mathrm{~cm}^3$
\\ \textit{Note:} The correct answer is derived from the equation for the volume of a gas under non-ideal conditions, using the formula \( V = Z \cdot \frac{nRT}{P} \), where \( Z \) is the compression factor (0.86), \( n \) is the number of moles (0.0082 mol), \( R \) is the ideal gas constant (0.0821 L·atm/(K·mol)), and \( P \) is the pressure (20 atm), which results in a calculated volume of 5.9 cm³.
\end{enumerate}

\section*{True / False}
\begin{enumerate}[label=\arabic*.]
\setcounter{enumi}{5}
\item \textbf{Answer:} True
\\ \textit{Options:} A) True; B) False
\\ \textit{Note:} The statement is true because the molarity is calculated by dividing the number of moles of sulfuric acid (approximately 0.102 moles from 10.0 grams) by the volume of the solution in liters (0.500 L), resulting in a concentration of 0.204 M.
\item \textbf{Answer:} True
\\ \textit{Options:} A) True; B) False
\\ \textit{Note:} Liquid hydrogen fluoride acts as a stronger and more selective electrophilic catalyst in the Friedel-Crafts alkylation reaction, promoting a more favorable reaction pathway for the formation of isopropylbenzene from propene and benzene compared to concentrated hydrobromic acid.
\item \textbf{Answer:} True
\\ \textit{Options:} A) True; B) False
\\ \textit{Note:} CsF is an ionic compound because it is formed by the transfer of electrons from cesium to fluorine, resulting in electrostatic attraction between the resulting ions, whereas PH 3 is a covalent compound due to the sharing of electrons between phosphorus and hydrogen atoms.
\item \textbf{Answer:} True
\\ \textit{Options:} A) True; B) False
\\ \textit{Note:} The antilogarithm base 10 of 0.8762 is 7.508, and subtracting 2 from this value results in 0.0752, confirming the statement as true.
\item \textbf{Answer:} False
\\ \textit{Options:} A) True; B) False
\\ \textit{Note:} The statement is false because the presence of ethanol, a solute, lowers the freezing point of water through colligative properties, meaning the freezing point of the solution will be below 0 degrees Celsius.
\end{enumerate}

\section*{Long Form Response}
\begin{enumerate}[label=\arabic*.]
\setcounter{enumi}{10}
\item \textbf{Answer:} To calculate the volume of the lunar atmosphere required to contain $10^6$ molecules and 1 millimole of gas at 100^\ensuremath{\circK}, we employ the ideal gas law, PV = nRT. Given that 1 mole of gas contains approximately 6.022 \ensuremath{\times} $10^23$ molecules, $10^6$ molecules correspond to about 1.66 \ensuremath{\times} 10^-17 moles. At a barometric pressure of 1.0 \ensuremath{\times} 10^-3 atm on the lunar surface and using R = 0.0821 L·atm/(K·mol), we find the volume for $10^6$ molecules to be approximately 1.04 \ensuremath{\times} 10^-4 L. For 1 millimole (0.001 moles), the same calculation yields a volume of approximately 6.24 \ensuremath{\times} $10^10$ L. Thus, the lunar atmosphere can contain these quantities of gas under the specified conditions.
\\ \textit{Note:} The answer correctly applies the ideal gas law (PV = nRT) to determine the volume of gas required at a specified temperature and pressure, accurately converting the number of molecules to moles and demonstrating a thorough understanding of the relationship between pressure, volume, and the amount of substance.
\item \textbf{Answer:} The equilibrium reaction I$_{2}$ + I- ⇌ I$_{3}$- illustrates the complexation of iodine in the presence of iodide ions, which significantly enhances the solubility of iodine in solution. The equilibrium constant K for this reaction can be expressed as K = [I$_{3}$-] / ([I$_{2}$][I-]), where the concentrations of the respective species are considered at equilibrium. Factors influencing this value include temperature, as changes can shift the equilibrium position, and the concentration of iodide ions, which directly affects the formation of the triiodide ion (I$_{3}$-). Additionally, solvent effects and the presence of other ions can alter the solubility and interaction dynamics, impacting the overall equilibrium constant. The increased solubility of iodine in KI solutions exemplifies Le Chatelier's principle, where the addition of I- drives the formation of I$_{3}$-, thus influencing K significantly, with a typical value around 330 at standard conditions.
\\ \textit{Note:} The answer accurately describes the complexation of iodine with iodide ions, defines the equilibrium constant in terms of the concentrations of the participating species, and identifies key factors such as temperature and iodide concentration that influence the equilibrium position, thereby providing a comprehensive analysis of the reaction.
\end{enumerate}

\vfill
\noindent\textit{End of Answer Key}
\end{document}